\section{Camera Follow}

Setelah player bergerak (Bab VI), kamera perlu mengikuti. Bab ini membahas pola camera follow. Kamera perlu mengikuti posisi bola agar pemain selalu melihat dari belakang-atas. Buat script \class{CameraController}:

\begin{unitycode}
using UnityEngine;

public class CameraController : MonoBehaviour
{
    public GameObject player;
    private Vector3 offset;

    void Start()
    {
        offset = transform.position - player.transform.position;
    }

    void LateUpdate()
    {
        transform.position = player.transform.position + offset;
    }
}
\end{unitycode}

\class{LateUpdate} dipanggil setelah semua Update. Dengan demikian, kamera bergerak setelah player bergerak di frame yang sama---menghindari jitter (getaran visual). Variabel \texttt{offset} menyimpan jarak awal antara kamera dan player; offset tersebut dijaga konstan agar sudut pandang tetap sama.

